Intensive-care records are irregularly sampled and mostly missing: in the PhysioNet/Computing in Cardiology Challenge 2012 cohort used here, only about one fifth of the hourly variable slots in the first two days contain a measurement (we report the exact fraction in Section~\ref{sec:setup}). Recurrent networks that treat missingness explicitly were introduced for exactly this setting. GRU-D~\cite{che2016recurrent} feeds an observation mask and a time-since-last-observation (``delta'') to a GRU and decays both the imputed inputs and the hidden state toward a default as the delta grows, and it has since served as the reference neural baseline for a long line of irregular-sampling models~\cite{shukla2019interpolation,horn2019set,rubanova2019latent,zhang2021graph,tipirneni2021self}. At the same time, large ICU benchmarks repeatedly find that logistic regression or gradient boosting on hand-built per-variable summaries is close to, or better than, sequence models~\cite{harutyunyan2017multitask,yche2021hirid,water2023yet}, in line with the general observation that tree ensembles remain hard to beat on medium-sized tabular data~\cite{grinsztajn2022why}.

This paper asks a deliberately narrow question on a small, open cohort (set-a of the PhysioNet 2012 challenge, \rhval{const/n_stays} stays): \emph{when all systems receive the same small tuning budget and are evaluated with repeated patient-level cross-validation, does a missingness-aware GRU beat a forward-filled GRU and beat summary-feature logistic regression and gradient boosting, and which ingredient of GRU-D (mask, delta, input decay, hidden decay) matters?} It is not a new method. Its value, if any, is a controlled comparison on a cohort small enough that effects are easily swamped by fold and seed noise, with every number traced to a logged run.

\textbf{Contributions.} (i) Five reimplemented systems (LR and gradient boosting on summaries, a forward-fill GRU, a GRU with mask and delta inputs, GRU-D) evaluated under one protocol: the same folds, preprocessing, metrics and tuning grid sizes, with \rhval{const/n_seeds} seeds of \rhval{const/n_folds}-fold patient-level CV and AUROC, AUPRC and Brier score. (ii) Ablations of the GRU-D decay mechanism and of the count features of the summary models, and an observation-horizon sweep. (iii) A hypothesis-by-hypothesis report, including that the headline recurrent model does not win.
